% !TeX root = ../概率论与数理统计笔记.tex
% 由主文档通过 \input 载入；请编译主文档。
\setcounter{chapter}{6} % 保留原书章号。
\chapter{条件期望与条件方差}
\label{chap:ps-conditional-expectation}

本章将“利用已知信息重新求平均”发展为系统工具。
先修内容为条件分布与数学期望；学习目标是理解条件期望本身可以是随机变量，并用它分解复杂计算。

\section{给定事件和随机变量的条件期望}
\label{sec:ps-conditional-expectation-definition}
由条件分布定义\emph{条件期望}（conditional expectation）。
区分固定取值 $y$ 下的数值 $\mathbb E[X\mid Y=y]$ 与随机变量 $\mathbb E[X\mid Y]$；以 $X$ 可积作为基本条件。
连续情形说明定义按条件分布的版本理解，不能把“给定取值”机械替换成正概率事件的比值。

\section{全期望公式与迭代性质}
\label{sec:ps-tower-property}
建立\emph{全期望公式}（law of total expectation）、线性性、已知量提出与信息嵌套下的迭代性质。
使用 $\sigma$ 代数表示信息的严格定义作为选读，明确可测性、积分刻画与几乎处处唯一性。
通过嵌套与非嵌套信息的比较说明塔式性质不能随意交换条件。

\section{条件方差与全方差公式}
\label{sec:ps-total-variance}
在 $X$ 二阶矩有限时定义\emph{条件方差}（conditional variance），建立\emph{全方差公式}（law of total variance）：
\begin{equation}
  \operatorname{Var}(X)
  =\mathbb E[\operatorname{Var}(X\mid Y)]
  +\operatorname{Var}(\mathbb E[X\mid Y]).
  \label{eq:ps-total-variance}
\end{equation}
解释组内波动与组间波动的含义，并通过中心化展开给出后续正文中的证明路线。

\section{条件期望作为最优平方损失预测}
\label{sec:ps-conditional-prediction}
在平方可积条件下，研究用 $Y$ 的可测函数预测 $X$ 的问题，证明条件期望使总体均方误差最小。
从误差正交分解解释“最优”的候选范围与几乎处处唯一性。
区分总体最优预测函数与由有限数据估计的回归函数，为第\ref{chap:ps-regression}章作铺垫。

\section{例题、习题与学习检查}
规划随机次数求和、分层总体均值与方差、给定联合密度求条件均值等例题。
随机和例题明确计数变量与被求和序列的独立性及所需矩条件。
习题重点检查“结果是数还是随机变量”、条件信息是否嵌套，以及平方损失结论的适用范围。






